\emph{(1).} This is the common-eigenvector lemma of the result filed alongside, synthesis.residual-s2-decouples-first-family-from-ckm-and-pmns (its clause (ii)); recalled here. Let $U_{u}$, $U_{d}$ be unitaries with $U_{q}^{\dagger}H_{q}U_{q}$ diagonal, the columns ordered by mass; $V = U_{u}^{\dagger}U_{d}$. With nondegenerate masses each column of $U_{q}$ is fixed up to a phase. If $x$ is a common unit eigenvector, it is (up to phase) column $i$ of $U_{u}$ and column $k$ of $U_{d}$ for some $i$, $k$, so $V_{ik} = e_{i}^{\mathsf{T}}U_{u}^{\dagger}U_{d}e_{k}$ has modulus one, and unitarity forces the rest of row $i$ and column $k$ to vanish. Also $[H_{u}, H_{d}]x = H_{u}\lambda_{d}x - H_{d}\lambda_{u}x = 0$, so $\det[H_{u}, H_{d}] = 0$; by Jarlskog's identity $\det[H_{u}, H_{d}] = 2I_{4}J\prod_{i<j}(m_{u_{j}}^{2} - m_{u_{i}}^{2})\prod_{i<j}(m_{d_{j}}^{2} - m_{d_{i}}^{2})$ (cited from memory) the invariant $J$ is zero. If $H_{u}$ and $H_{d}$ commute they are simultaneously diagonal and every column of $U_{d}$ is a column of $U_{u}$ up to phase.

\emph{(2).} In Fact \claimref{standard-model.imported-clebsch-structure-seesaw} each coupling enters $M_{u}$ and $M_{d}$ multiplied by complex vevs. If every coupling is $\mu_{a}Y$ for one matrix $Y$, then $M_{u} = (\sum_{a}\mu_{a}v^{a}_{u})Y$ and $M_{d} = (\sum_{a}\mu_{a}v^{a}_{d})Y$, so $M_{u} = rM_{d}$ with $r$ the ratio (when $M_{d} \ne 0$). Then $H_{u} = |r|^{2}H_{d}$: the eigenbases coincide, $V = 1$, and the singular values satisfy $m_{u_{i}} = |r|m_{d_{i}}$. This covers the one rung coupling of Corollary \claimref{standard-model.one-vev-vector-relations-breaking-pattern} under the $c^{*}$ reading, where $r = 1$, and the complex bidoublet of Proposition \claimref{standard-model.bidoublet-complex} when its two doublets carry the same family matrix, where $|r|$ is the ratio of the two vevs (a $\tan\beta$). The check verifies $H_{u} = |v_{u}/v_{d}|^{2}H_{d}$ and $[H_{u}, H_{d}] = 0$ for a generic complex $Y$ and three pairs of vevs, and, as a control, that two independent family matrices give $\det[H_{u}, H_{d}] \ne 0$.

\emph{(3).} By Lemma \claimref{neutrino.breaking-relation} and Proposition \claimref{synthesis.residual-s2-fixes-120-family-form}(i), the $S_{3}$-invariant symmetric family matrices are spanned by $1$ and $\J$ (equivalently $\Pi$ and $\J$), and no nonzero antisymmetric matrix is $S_{3}$-invariant, so no $120$ enters. Since $\J^{2} = 3\J$, this span is closed under products and adjoints (for complex coefficients), so $M_{u}$, $M_{d}$, $H_{u}$, $H_{d}$ all lie in it. $\J$ has eigenvalue $3$ on $u' = (1,1,1)$ and $0$ on its orthogonal plane, so $\alpha 1 + \beta\J$ has eigenvalues $\alpha$ (twice) and $\alpha + 3\beta$, and $H_{q}$ has $|\alpha_{q}|^{2}$ twice: two families of each sector are degenerate. Clause (1) then does not apply (it assumes nondegenerate masses) and the mixing is undetermined; the record's nondegenerate masses exclude this case outright. The check verifies the spectra for a complex instance, $H_{u}$ with $5, 5, 290$ and $H_{d}$ with $9, 9, 90$. The general residual-$S_{2}$ case, of which this is a special case, is the result filed alongside, synthesis.residual-s2-decouples-first-family-from-ckm-and-pmns.

\emph{(4).} Record values (cited from memory; running masses at $M_{Z}$ in MeV, from the table of Xing, Zhang and Zhou, Phys. Rev. D 77 (2008) 113016): $m_{u} \approx 1.27$, $m_{d} \approx 2.90$, $m_{c} \approx 619$, $m_{s} \approx 55$, $m_{t} \approx 1.717 \times 10^{5}$, $m_{b} \approx 2.89 \times 10^{3}$, so $m_{u}/m_{d} \approx 0.44$, $m_{c}/m_{s} \approx 11$, $m_{t}/m_{b} \approx 59$; $|V_{us}| \approx 0.225$ and $J \approx 3 \times 10^{-5}$ (Particle Data Group). Three different ratios contradict (2) at $M_{Z}$. To carry the conclusion to $M_{I}$: QCD running multiplies the up and down masses of each family by the same factor at one loop, and the electroweak and Yukawa corrections change the ratios by factors of order one, not by the factor of about $130$ between $m_{t}/m_{b}$ and $m_{u}/m_{d}$; and the one-loop running of the Yukawa matrices preserves $[H_{u}, H_{d}] = 0$ (the flow of $H_{u}$ and $H_{d}$ is generated by polynomials in them), so a vanishing mixing angle at $M_{I}$ stays zero at $M_{Z}$ (both cited from memory). So $V \ne 1$ at $M_{Z}$ rules out (2) and (3) as the source of the observed spectrum. Under the $c^{*}$ reading the second coupling must therefore be a family matrix not proportional to the first. For three families with $J \ne 0$ it must give $\det[H_{u}, H_{d}] \ne 0$. The stored record-level fits, whose complex couplings reproduce the CKM numbers, are of this kind.
